\subsection{Exact decomposition}

Let \(R_t\) denote average remuneration in period \(t\). The population and
remuneration concept are defined separately in each paper. Educational
attainment partitions that population into mutually exclusive groups indexed by
\(g\):
\begin{equation}
  R_t = \sum_g s_{g,t} r_{g,t},
  \label{eq:remuneration_identity}
\end{equation}
where \(s_{g,t}\) is group \(g\)'s share of the relevant population and
\(r_{g,t}\) is its average remuneration.

For two periods, the change in average remuneration can be written as
\begin{equation}
  R_1-R_0 =
  \sum_g \bar{s}_g\left(r_{g,1}-r_{g,0}\right)
  +
  \sum_g \bar{r}_g\left(s_{g,1}-s_{g,0}\right),
  \label{eq:symmetric_decomposition}
\end{equation}
where
\(\bar{s}_g=(s_{g,1}+s_{g,0})/2\) and
\(\bar{r}_g=(r_{g,1}+r_{g,0})/2\).
The first term measures changes in remuneration within educational groups. The
second measures changes in the educational composition of the relevant
population. This symmetric decomposition is exactly additive and does not
depend on choosing either period as the reference.

Equation~\eqref{eq:symmetric_decomposition} is an accounting identity. The
within-group component is not a measure of productivity, and the composition
component does not identify the causal effect of education. Changes in either
component may reflect selection, labor supply, labor demand, institutions,
hours, occupations, sectors, or survey coverage.
